\solchapter{qinfo}

\sol{qinfo}{11} List the four possible states (no flip, or a flip on qubit 1, 2, or 3) and compute $Z_1Z_2$ and $Z_2Z_3$ for each. Do $\ket{000}$ and $\ket{111}$ (and their flipped versions) give the same values?

\sol{qinfo}{12} After the box, the first qubit is $[(-1)^{f(0)}\ket0 + (-1)^{f(1)}\ket1]/\sqrt2$. What is this when $f(0) = f(1)$, and when $f(0)\ne f(1)$?

\sol{qinfo}{13} Group $C = A(B+B') + A'(B-B')$ before squaring. You will need $(B+B')(B-B') = -[B,B']$.

\solchapter{chem}

\sol{chem}{11} Expand $\sqrt{1\pm\kappa}$ to second order in $\kappa$. The first-order terms cancel in the sum.

\sol{chem}{12} Problem~5 told you which orbital moves more. For O$_2$ and N$_2$, fill the valence molecular orbitals in the order $\sigma_{2s},\sigma_{2s}^*,\pi_{2p},\sigma_{2p},\pi_{2p}^*,\sigma_{2p}^*$ (for O$_2$ the $\sigma_{2p}$ drops below the $\pi_{2p}$, but the count is the same).

\solchapter{nuclear}

\sol{nuclear}{10} The half-life is the inverse of (attempt rate)$\times$(tunnelling probability). Write $v$ in terms of $E$; then evaluate $2\pi\eta = 4\pi Z\alpha c/v$ at the two energies.

\sol{nuclear}{11} Write a diffusion equation with a source: $\partial n/\partial t = D\nabla^2n + (\nu-1)n/\tau$, with $D = \ell v/3$ and $\tau = \ell/v$. The slowest-decaying mode in a sphere has $\nabla^2n = -(\pi/R)^2n$.

\sol{nuclear}{12} The Coulomb barrier is $Z_1Z_2\times1.44\un{MeV\,fm}/R$, with $R$ the sum of the radii. For the s-process, follow what happens to bismuth-209 after it captures a neutron.

\solchapter{bio}

\sol{bio}{10} How long does a molecule stay near the receptor before diffusing away? In that time, how many are there, and how many such independent looks fit in $T$?

\sol{bio}{11} Write the rate of forward steps (the gas, at $T_1$, supplies $\epsilon$ plus the work) and of backward steps (the pawl, at $T_2$, is lifted by $\epsilon$). Compare. Then find the heat and work per step.

\sol{bio}{12} After activation, the complex either proceeds (rate $W$) or releases the substrate (rate $\ell'$). If $W\ll\ell'$, what fraction proceeds, and how does it differ for right and wrong substrates? For the second part, ask what equilibrium alone can achieve.

\solchapter{medical}

\sol{medical}{11} For backprojection, write $\delta(x\cos\theta+y\sin\theta)$ in polar coordinates and integrate over $\theta$. For the filter, write the inverse two-dimensional Fourier transform in polar coordinates and look at the Jacobian.

\sol{medical}{12} Write the absorbance at each wavelength as $cL[s\epsilon_{\text{ox}} + (1-s)\epsilon_{\text{de}}]$ and take the ratio: what drops out? Then list everything in the light path that does not change with the heartbeat.

\solchapter{geo}

\sol{geo}{10} Draw the chord from the earthquake to the edge of the shadow, $104^\circ$ around the Earth. How close does it come to the centre?

\sol{geo}{11} The Moon's orbital angular momentum is $M_m\sqrt{GM_ea}$. Differentiate with respect to time, set the Earth's loss of spin angular momentum equal to the Moon's gain, and convert $\dd\omega/\dd t$ into a change of the day's length.

\sol{geo}{12} Find the equilibria on each of the three pieces of the absorbed curve. For stability, ask whether a small warming makes the planet emit more than it absorbs.
